\documentclass[10pt]{article}
\usepackage[T1]{fontenc}
\usepackage[utf8]{inputenc}
\usepackage[french]{babel}
\usepackage{amsmath,amssymb,mathrsfs,mathpazo}
\usepackage{geometry}
\geometry{a4paper,margin=2cm}
\usepackage{graphicx,array,tabularx,tabularray,colortbl,multirow,multicol,booktabs,enumitem,pifont,ulem,fancybox,diagbox,cancel}
\usepackage{pstricks,pstricks-add,pst-plot,pst-tree,pst-node,pst-text,pst-func,pst-math,pst-bspline,pst-3dplot,pst-eucl,pst-solides3d}
\let\pscalculate\undefined
\usepackage{pst-bezier}
\usepackage{tikz,pgfplots,tkz-tab}
\usetikzlibrary{arrows,arrows.meta,patterns,calc,positioning,trees,decorations.markings,shapes.geometric,shapes.misc,shapes.arrows,intersections,decorations.pathmorphing,decorations.pathreplacing,backgrounds,shadows,babel}
\pgfplotsset{compat=1.18}
\usepgfplotslibrary{fillbetween}
\usepackage[np]{numprint}
\usepackage{listings,makecell,pas-tableur,twoopt}
\let\barre\undefined
\usepackage[pstricks]{tablvar}
\definecolor{darkgreen}{RGB}{0,100,0}
\providecommand{\red}{\color{red}}
\providecommand{\blue}{\color{blue}}
\providecommand{\green}{\color{green}}
\providecommand{\black}{\color{black}}
\providecommand{\white}{\color{white}}
\providecommand{\gray}{\color{gray}}
\providecommand{\cyan}{\color{cyan}}
\providecommand{\magenta}{\color{magenta}}
\providecommand{\yellow}{\color{yellow}}
\providecommand{\violet}{\color{violet}}
\providecommand{\orange}{\color{orange}}
\providecommand{\euro}{\texteuro}
\providecommand{\EUR}{\texteuro}
\providecommand{\ieme}{\textsuperscript{e}}
\providecommand{\degres}{\ensuremath{^\circ}}
\providecommand{\degree}{\ensuremath{^\circ}}
\providecommand{\vv}[1]{\overrightarrow{#1}}
\providecommand{\Coord}[1]{#1}
\providecommand{\decosix}{}
\providecommand{\starredbullet}{\ensuremath{\bullet}}
\pagestyle{empty}
\setlength{\parindent}{0pt}
\usepackage{adjustbox}\DeclareUnicodeCharacter{2236}{:}\DeclareUnicodeCharacter{2190}{\ensuremath{\leftarrow}}
\providecommand{\R}{\mathbb{R}}
\providecommand{\N}{\mathbb{N}}
\providecommand{\Z}{\mathbb{Z}}
\providecommand{\Q}{\mathbb{Q}}
\providecommand{\D}{\mathbb{D}}
\providecommand{\C}{\mathbb{C}}
\providecommand{\e}{{\mathrm{e}}}
\providecommand{\vect}[1]{\overrightarrow{#1}}
\providecommand{\vectt}[1]{\overrightarrow{\mathrm{#1}}}
\providecommand{\vvt}[1]{\overrightarrow{\mathrm{#1}}}
\providecommand{\np}[2][]{{#2~#1}}
\providecommand{\vv}[1]{{\overrightarrow{#1}}}
\renewcommand{\d}{{\mathrm{d}}}
\providecommand{\degres}{{^\circ}}
\providecommand{\degre}{{^\circ}}
\providecommand{\blue}{\color{blue}}
\providecommand{\red}{\color{red}}
\providecommand{\enskip}{\quad}
\providecommand{\llbracket}{\lbrack\!\lbrack}
\providecommand{\rrbracket}{\rbrack\!\rbrack}
\providecommand{\euro}{{\text{€}}}
\providecommand{\no}{{\text{n°}}}
\providecommand{\interval}[2]{{[#1;#2]}}
\providecommand{\intervFO}[2]{{[#1;#2[}}
\providecommand{\dotfill}{\ldots}

\providecommand{\ds}{\displaystyle}
\providecommand{\textding}[1]{#1}
\providecommand{\Oij}{\(\mathrm{O};\vec{\imath},\vec{\jmath}\)}
\providecommand{\Oijk}{\(\mathrm{O};\vec{\imath},\vec{\jmath},\vec{k}\)}
\providecommand{\Ouv}{\(\mathrm{O};\vec{u},\vec{v}\)}
\providecommand{\intervFF}[2]{\ensuremath{[#1;#2]}}
\renewcommand{\euro}{\ifmmode\text{\texteuro}\else\texteuro\fi}
\DeclareUnicodeCharacter{00B0}{\ifmmode^{\circ}\else\textdegree\fi}
\DeclareUnicodeCharacter{20AC}{\ifmmode\text{\texteuro}\else\texteuro\fi}
% Énoncés : Éducation nationale. Transcription et corrigés : APMEP, Denis Vergès et contributeurs. Découpage et classement : CoopMaths, Sébastien Lozano, Éric Elter et contributeurs. Adaptation de la présentation : Maths Explorer.
% CC BY-SA : https://coopmaths.fr/www/about/
% https://coopmaths.fr/alea/static/e3c/2020/tex/e3c_2020_00_sujet21_3_cor.tex
\begin{document}
\noindent\textit{APMEP / CoopMaths. CC BY-SA. Présentation adaptée par Maths Explorer.}\par\medskip
\subsection*{1.}
	
	\(\overrightarrow{AB} \left( \begin{array}{c} -1 \\ 5 \end{array} \right)\) et \(\overrightarrow{AC} \left( \begin{array}{c} -5 \\ 2 \end{array} \right)\), d’où \(\overrightarrow{AB} \cdot \overrightarrow{AC} = 5 + 10 = 15\).
	
	\subsection*{2.}
	
	a. Soit \(\overrightarrow{AB} \cdot \overrightarrow{AD} = \overrightarrow{AB} \cdot (\overrightarrow{AC} + \overrightarrow{CD}) = \overrightarrow{AB} \cdot \overrightarrow{AC}\), car \(\overrightarrow{AB}\) et \(\overrightarrow{CD}\) sont orthogonaux.
	
	b. \(\overrightarrow{AB} \cdot \overrightarrow{AD} = \overrightarrow{AB} \cdot \overrightarrow{AC}\) soit \(AB \times AD = 15\).
	
	Or
	\[
	AB^2 = 1 + 25 = 26, \quad \text{donc} \quad AB = \sqrt{26}, \quad \text{donc} \quad AD = \dfrac{15}{\sqrt{26}}.
	\]
	
	\subsection*{3.}
	
	Dans le triangle \(ACD\) rectangle en \(D\), le théorème de Pythagore s’écrit :\\
$
	CD^2 + AD^2 = AC^2$  soit $ CD^2 + \left(\dfrac{15}{\sqrt{26}}\right)^2 = 5^2 + 2^2 = 29$, \\
	 d’où $ AD^2 = 29 - \dfrac{225}{26} = \dfrac{29 \times 26 - 225}{26} = \dfrac{754 - 225}{26} = \dfrac{529}{26}.$
	
	Donc la hauteur du triangle \(ABC\) issue de \(C\) a pour longueur
	\[
	CD = \sqrt{\dfrac{529}{26}}.
	\]
	
	\subsection*{4.}
	
	En prenant comme base \([AB]\), avec \(AB = \sqrt{26}\) et comme hauteur \([CD]\), avec \(CD = \sqrt{\dfrac{529}{26}}\).
	
	\[
	A(ABC) = \dfrac{AB \times CD}{2} = \dfrac{\sqrt{26} \times \sqrt{\dfrac{529}{26}}}{2} = \dfrac{\sqrt{529}}{2} = \dfrac{23}{2} = 11,5.
	\]

\end{document}
